1) Initial sticker values (only what matters for the front-face sum)
- The 8 non-center stickers on the original Front face (F) are all 1.
- The 8 non-center stickers on the original Back face (B) are all 1.
- Every non-center sticker on U, D, L, R is 0.
- All 6 centers are 2.
Therefore, after any moves:
- The front-center sticker is still the original front center, hence it is always 2.
- The sum on the front face = 2 + (number of "1-stickers" currently on the front face).

2) Reduce the problem
We only need to count how many of the sixteen 1-stickers (the 8 perimeter stickers on F and the 8 perimeter stickers on B) end up facing Front (i.e. lying on the front face) after the move sequence

$R, U, B, R, L, L, R, L, R, L$.

3) Model for tracking stickers
Use coordinates with axes: x right, y up, z front. Faces are:
- Front F: normal (0,0,1), Back B: (0,0,-1), Right R: (1,0,0), Left L: (-1,0,0), Up U: (0,1,0), Down D: (0,-1,0).

Track each 1-sticker as a pair (position p, outward normal n). A sticker is on the front face exactly when its final normal is (0,0,1).

Rotation rules used (standard cube notation):
- R: rotate layer x=+1 by +90 degrees about x: $(x,y,z)\mapsto (x,-z,y)$ and $n\mapsto (n_x,-n_z,n_y)$.
- L: rotate layer x=-1 by +90 degrees about x: same formula but applied only when x=-1.
- U: rotate layer y=+1 by +90 degrees about y: $(x,y,z)\mapsto (z,y,-x)$ and $n\mapsto (n_z,n_y,-n_x)$.
- B: rotate layer z=-1 by -90 degrees about z: $(x,y,z)\mapsto (y,-x,z)$ and $n\mapsto (n_y,-n_x,n_z)$.

4) Track the front-perimeter 1-stickers (those initially on F)
Denote the five front stickers that end up on the front face at the end (computations follow the above rotation rules, step-by-step through the 10 moves):

(a) FD (front-down edge) starts at p=(0,-1,1), n=F.
None of the moves R,U,B,L affects this sticker's cubie-layer (it is never in y=+1, x=+-1, z=-1), so it stays on the front face.
=> FD ends with n=F.

(b) FR (front-right edge) starts at p=(1,0,1), n=F.
Applying the 10 moves gives final normal n=F.
=> FR ends on the front face.

(c) FL (front-left edge) starts at p=(-1,0,1), n=F.
Applying the 10 moves gives final normal n=F.
=> FL ends on the front face.

(d) FUR (front-up right corner sticker) starts at p=(1,1,1), n=F.
Applying the 10 moves gives final normal n=F.
=> FUR ends on the front face.

(e) FDL (front-down left corner sticker) starts at p=(-1,-1,1), n=F.
Applying the 10 moves gives final normal n=F.
=> FDL ends on the front face.

The remaining three original front 1-stickers (FU, FUL, FDR) end with normals different from F (they end on side faces), hence they do NOT contribute to the front-face count.

5) Track the back-perimeter 1-stickers (those initially on B)
Applying the same update rules to the 8 back 1-stickers (BU, BR, BD, BL and the four back corners) shows that none finishes with normal F.
=> 0 back-1-stickers contribute to the final front face.

6) Count 1-stickers on the final front face
From Steps 4-5:
- Number of 1-stickers on the final front face = 5.

7) Compute the required sum on the front face
- Front center = 2.
- 5 perimeter stickers are 1.
- The remaining 3 perimeter stickers are 0.

Sum = 2 + 5*1 + 3*0 = 7.

FINAL ANSWER: 7
